% ============================================================
% 2. RELATED WORK
% ============================================================
\section{Related Work}

Our work intersects three research areas: persona-consistent dialogue generation, self-evolving agent architectures, and the side effects of RLHF alignment on open-ended generation.

\subsection{Persona-Consistent Dialogue}

Maintaining a consistent character identity across multi-turn conversations has been studied extensively. Early approaches encode persona as fixed textual descriptions---PersonaChat \cite{zhang2018personalizing} introduced persona-conditioned dialogue where models receive a set of personality sentences. Subsequent work improved consistency through mutual persona perception \cite{liu2020you}, persona-augmented knowledge \cite{majumder2020like}, and contrastive learning of persona representations \cite{cao2022model}.

More recent systems scale persona control to LLMs. CharacterGLM \cite{zhou2023characterglm} fine-tunes models on large-scale character dialogue datasets. LIGHT \cite{urbanek2019learning} situates characters in grounded fantasy environments. RoleLLM \cite{wang2023rolellm} uses role-profile extraction and context-based instruction tuning. CharacterLLM \cite{shao2023character} trains agents on historical character biographies.

These approaches focus on \emph{who} the character is (knowledge, identity, backstory) but largely treat emotional expression as a byproduct of persona grounding. They do not model \emph{how} emotional state evolves dynamically across turns---an agent may be consistent in personality but monotonous in expression. TPE complements these methods by operating on the behavioral signal layer rather than the identity layer.

\subsection{Self-Evolving Agent Architectures}

Generative Agents \cite{park2023generative} introduced the paradigm of long-term memory, reflection, and planning for believable character simulation. Their reflection mechanism---periodically summarizing memories into higher-level observations---directly inspires Reflexion \cite{shinn2023reflexion}, which applies similar self-correction to task-solving: agents reflect on failures and adjust strategies through a verbal reinforcement loop.

Promptbreeder \cite{fernando2023promptbreeder} takes self-evolution further by mutating both task prompts and the mutation operators themselves. CAMEL \cite{li2023camel} explores multi-agent role-playing through inception prompting. MemGPT \cite{packer2024memgpt} introduces hierarchical memory management for LLMs, enabling virtual context management beyond the model's native window.

Our system differs from these architectures in two key respects. First, TPE's memory is organized by \emph{behavioral similarity} in continuous signal space, not by semantic content or temporal recency. This means retrieval is driven by ``how did I behave in emotionally similar situations?'' rather than ``what happened recently?'' Second, our experimental results (\S\ref{sec:reflexion}) demonstrate that Reflexion-style self-correction, while effective for task convergence, actively degrades persona consistency---a finding that questions the universal applicability of reflective self-correction across agent types.

\subsection{RLHF Side Effects on Open-Ended Generation}

RLHF \cite{ouyang2022training} and its variants (DPO, RLAIF) have become the standard post-training paradigm for instruction-following LLMs. While highly effective for safety and helpfulness, alignment training introduces well-documented side effects.

Mode collapse is a recognized concern: aligned models converge toward a narrow distribution of ``helpful assistant'' responses \cite{kirk2023understanding}. Sycophancy---aligned models tending to agree with users rather than maintain consistent positions---has been empirically demonstrated \cite{perez2022discovering}. The ``alignment tax'' \cite{askell2021general} describes capability degradation from overly conservative training.

In role-playing contexts, these effects manifest as what we term ``RLHF collapse''---the agent abandons its designated persona to produce safe, empathetic, counselor-style responses. Our results on gpt-5-mini (\S\ref{sec:alignment}) provide empirical evidence of this phenomenon: a heavily aligned model systematically overrides persona instructions with template therapeutic language, achieving Collapse Index scores $22\times$ higher than less-aligned models.

TPE combines behavioral context modulation (few-shot memory examples shaped by gravitational retrieval) with prompt-level persona instruction. The behavioral context provides an additional control channel beyond prompt engineering alone: even when the prompt-level instruction is partially overridden by alignment, the few-shot examples still shape the distribution of generated text. However, this dual-channel approach is not sufficient to overcome strong alignment, as the gpt-5-mini failure demonstrates (\S\ref{sec:alignment}).
